\documentclass[11pt,a4paper]{article}
\usepackage{amsmath,graphicx,psfrag,pstcol,amsfonts,stmaryrd,amssymb}
\usepackage{hyperref}

\usepackage{enumerate} % Custom item numbers for enumerations

\usepackage[ruled]{algorithm2e} % Algorithms

\usepackage[framemethod=tikz]{mdframed} % Allows defining custom boxed/framed environments

% --- FONT SETUP FOR PDFLATEX ---
\usepackage[T1]{fontenc}
\usepackage[utf8]{inputenc}

% 'osf' = Old Style Figures (123 instead of 123)
% 'sc'  = Real Small Caps support
\usepackage[sc,osf]{mathpazo} 

% Palatino generally looks better with slightly wider line spacing
\linespread{1.05}

%more spacing between paras
\usepackage{parskip}

% Recommended for high-quality text in pdflatex
\usepackage{microtype} 
% Standard margins
\usepackage[margin=1in]{geometry}

%remove itemize spacing
\usepackage{enumitem}

% For URLs
\usepackage{hyperref}
% Make links dark blue instead of ugly boxes
\hypersetup{
    colorlinks=true,
    linkcolor=blue,
    filecolor=magenta,
    urlcolor=blue,
}

% Reduce title spacing
\usepackage{titling}
\setlength{\droptitle}{-4em}

%----------------------------------------------------------------------------------------
%	ASSIGNMENT INFORMATION
%----------------------------------------------------------------------------------------

\title{ 3F7 Information Theory and Coding \\ \textit{Full Technical Report}} % Title of the assignment

\author{Liu Zihe \\ zl559} % Author name and email address

\date{University of Cambridge --- 6 December 2025} % University, school and/or department name(s) and a date

%----------------------------------------------------------------------------------------

\begin{document}

\maketitle % Print the title

\section{Introduction}

Brief intro of whole Report

\section{Overview of Short Lab}

summary of the lessons learned in the short lab (this should be fairly short, e.g., one
page, and stick to the essentials)

refer to \texttt{3F7/3F7lab.ipynb} for this

\section{Coding for Discrete Stationary Sources}

\subsection{Background: Entropy of Stochastic Processes}

Explain sources with memory, stationary sources etc

Explain core math behind section 1.1 in the handout. Refer to cover and thomas elements of information theory chapter 4 (should be attached and fed into your context window). Include stationarity and markov chains and their associated entropy math.

Note to self: memoryless means iid. Discrete stationary process can (and usually does have memory). Conditional probabilities depend only on relative positions (shift invariant). We are using strict stationarity here btw.

\subsubsection{Theorem 1}

proofs for all 4 statements in theorem1. For statement4, take the simplest proof. I have provided some context from cover and thomas as reference; you don't have to follow their method

\subsubsection{Theorem 2}

\subsection{Coding Schemes for Discrete Stationary Processes}

think
about what this implies in terms of the coding techniques you’ve learned during the short lab.

\textit{How would you implement Huffman and arithmetic coding for discrete stationary sources to
    achieve a good compression rate close to the entropy rate of a discrete stationary source?}

\textit{Can
    you think of a practical compression technique for which n can be made very large and hence
    have performance approaching the entropy rate, and how would this work in practice?}

\subsection{Entropy Rates of Discrete Stationary Processes}

\textit{For measuring the entropy rate, take the text file hamlet.txt and try to measure the
    entropy resulting from estimating the probabilities of couples, triples, quadruples etc. of consecutive symbols in the text file. We’ve provided a Jupyter notebook to help you measure these}

You will observe that while your estimates of H1, H2 and possibly H3 are fairly
accurate, the estimated entropies keep decreasing with apparently no positive lower bound,
implying possibly that H-inf = 0.

\textit{Do you really believe that the entropy rate of text is zero? What do you think may be going wrong in this experiment?}

\section{Adaptive Coding}

\subsection{Background: Universal Codes}

Using section 13 of Cover and Thomas, explain the math background behind this stuff:

The problem of probability estimation is even more interesting when you consider that
a good compression algorithm will in fact not need a separate stored probability model.
Compression can proceed starting from a startup probability model that assigns for example
a uniform distribution to all symbols, and adapt its model “on the fly” as you progress through
the file. As long as you always encode a symbol using the “old” model before updating your
estimated probabilities to get a “new” model, the decoder will be able to decode the source
symbol and update the source model to mirror exactly the updates done by the encoder.

Prove this estimator stuff:

In IB Paper 7, you’ve learned about the Beta distribution (Lecture slides 3) for estimat-
ing the unknown parameter of a source of Bernoulli random variables. If p is the unknown
Bernoulli parameter p = PX (1) of a binary source, by assuming that p is a random variable P
uniformly distributed between zero and one (“uniform prior”), then the distribution of P con-
ditioned on the data if you’ve observed n0 zeros and n1 ones is a Beta distribution Beta(n0, n1).
The Laplace estimator is the expected value of the Beta distribution E[P |X1, . . . , Xn] where
X1, . . . , Xn is the data seen after n = n0 + n1 observations
E[P |X1, . . . , Xk] = n1 + 1
n0 + n1 + 2 ,
which is the same as saying that before you start observing source outputs, you assume that
you’ve already seen one 0 and one 1. You can prove all of the above if you feel adventurous.

\subsection{Adaptive Arithmetic Coding}

Without actually implementing an adaptive arithmetic code, you can get a good estimate
of its performance by assuming that the length of the code sequence for a string of data
x1, x2, . . . , xn whose estimated probability is

Based on this estimate, you can investigate and answer the following questions by conducting
a few statistical measurements on the file hamlet.txt:

\subsubsection{Effect of length $n$ on performance}

how does the performance of adaptive arithmetic coding approach the entropy when
probabilities are adapted at each step using the estimator in (3) in function of the data
length n? You could plot L/n calculating L based on (4) using the first n symbols
in Hamlet, for n going from one to, say, 10,000, using, say delta = 1 (as in the Laplace
estimator). Does the resulting performance tend to the entropy of the distribution?


\subsubsection{Effect of bias $\delta$ on performance}

How does the bias delta affect the performance of the compression algorithm? Try various
values for delta and give an indication of how delta is best picked. If you’re feeling adventurous,
you could try to search for the optimal delta for a few compressed files using a Python
search library.

\subsubsection{Adaptive and non-adaptive coding schemes}

Of course, there may be some scenarios where storage of
such a model can indeed be discounted, for example if compressing many files with a unique
probability model that has been measured offline and can be viewed as “part of the software
package” rather than stored specifically for each file as we have done in the short lab.

For
other scenarios, you could also compare the performance of the two while accounting for the
fixed extra storage need for the offline model.

What are your conclusions from this compari-
son?

Is there a “sweet spot” for the file length from when it becomes more advantageuous to
store the model offline than to measure it on the fly?

\section{Adaptive Coding for Discrete Stationary Processes}

How does adaptive probability modelling apply to sources with mem-
ory?

What are the advantages and drawbacks of considering ever longer block lengths or
context lengths, if you are working with estimated probabilities rather than assuming that
probabilities are just given to you and accurate as we assumed in the previous section?

\section{Conclusion}

Some guidelines for you (the LLM):

\begin{itemize}[nosep]
    \item Please ignore the section on rate distortion theory, this is wholly unneccessary.
    \item Whenever you reference proofs, *please* add a reference the relevant section of the guide (e.g. cover and thomas)
    \item Please do not supply any python code to this FTR report, unless it is core pseudocode (please notify me if you plan to do this)
\end{itemize}


\end{document}

